% Section 7 — Empirical anchoring
\section{Empirical anchoring}
\label{sec:anchoring}

The preceding sections are self-contained mathematics: \cref{thm:irrelevance}
closes the information channel at the optimum, \cref{thm:bounded-gap} bounds
what remains for a bounded learner, and \cref{sec:parity-defect} measures the
residual when parity is only operational. This section connects those results
to the three frozen outcomes of \texttt{preprint-v1} --- the T1 clean null, the
T3 ceiling, and the T2 moderate-but-absorbed positive --- by naming, for each,
the theoretical regime it falls in. The exercise is strictly
\emph{retrodiction}: every empirical number here is read off the frozen record,
none is produced by the theory, and no new empirical claim is built. The point
is to show that the recorded outcomes are the ones the theory
\emph{already permits}, and to be exact about where the mapping from an
idealized learner to a single frozen LLM stops being a deduction and becomes an
analogy.

Three caveats bind the whole section and are stated once here rather than
repeated. First, the frozen artifact is one fixed model queried under five
representations, not an ERM ranging over a hypothesis class with a tunable
sample size $\nsample$; the PAC/Rademacher objects of \cref{sec:bounded} are
therefore read onto it by analogy, with the model's in-context behaviour
standing in for $\widehat f_{\rep}$ and the task corpus for the draw $S$
(disclosed again at each use). Second, the control is C1${}^{+}$ at
\emph{operational} parity (\cref{def:parity}), not absolute parity, so every
invocation of \cref{thm:irrelevance} below inherits idealization~(ii) of
\cref{rem:slogan} and is corrected by the \cref{sec:parity-defect} defect.
Third, none of what follows is a test: the theory is consistent with the record,
which is necessary but not sufficient for the theory to be true.

\subsection{The T1 null as the predicted information-limit outcome}
\label{subsec:anchor-t1}

The central H1 test on T1 --- the cleanest unconfounded cell, with no rubric
ceiling --- recorded a rank-biserial effect $r_{\mathrm{rb}}=+0.10$ at raw
$p=0.36$, a clean non-rejection of C4${}\le{}$C1${}^{+}$. Under the theory this
is not a weak or disappointing positive; it is the \emph{predicted}
outcome of the information-limit regime, and that prediction is exactly what
the honesty discipline forbids us to soften.

The mechanism is \cref{thm:irrelevance}. C4 and C1${}^{+}$ are constructed for
operational information parity, $\ladder{4}\parityeq\ladderplus$
(\cref{tab:parity-ledger}), so their information gaps vanish,
$\infogap{\ladder{4}}=\infogap{\ladderplus}=0$, and the at-optimum risks
coincide: $\bayesrisk(\ladder{4})=\bayesrisk(\ladderplus)$. At the information
limit there is simply no quantity for a format effect to be made of. What the
empirical comparison then measures is the residual of
\cref{thm:bounded-gap}: the advantage $\advantage(\ladderplus,\ladder{4})$ a
\emph{bounded} predictor can extract once the information channel is closed,
which the theorem caps at the capacity budget
$\Phi(\ladderplus,\ladder{4};\nsample,\delta)$ --- a quantity built from the
approximation gap and the induced-class Rademacher complexity, containing no
mutual information. A near-zero observed effect is the statement that this
budget was small (or unrealized) in T1's regime.

\cref{prop:regime} says when that is expected. A near-zero advantage is the
generic reading of its non-activation clause~(ii): absent a genuine capacity
separation --- a positive approximation gap $\approxerr{\ladder{4}}<
\approxerr{\ladderplus}$, or a finite-$\nsample$ estimation term comparable to
it --- the permitted window is empty and the bound itself sits near zero. A
frontier model on single-function, non-recursive Python is plausibly on the
benign side of clause~(i)'s threshold (it generalizes from in-context structure
rather than failing to), so the regime that \emph{could} have produced a large
parity-preserving effect is the one the theory marks as requiring a separation
T1 evidently did not supply. The null is thus doubly expected: forced toward
zero at the optimum by \cref{thm:irrelevance}, and not pushed away from zero by
the bounded term because T1 exhibits no capacity separation the theory can point
to.

\idealization{This subsection leans on both idealizations of
\cref{rem:slogan}. (i)~The exact zero is the \emph{Bayes-optimal} risk equality;
the frozen LLM is not the posterior $\post{\rep}$, so what is observed is the
bounded residual of \cref{thm:bounded-gap}, and the claim is only that this
residual is small, not that it is structurally zero. (ii)~Parity is operational,
so $\infogap{\cdot}=0$ holds up to the \cref{sec:parity-defect} defect rather
than exactly; the C1${}^{+}$/C4 byte asymmetry of \cref{prop:abs-strict} is the
admitted slack. The mapping of $\widehat f_{\rep}$ onto a single model's
in-context predictions is the analogy flagged in the preamble.}

\subsection{The T3 ceiling as sub-frontier saturation}
\label{subsec:anchor-t3}

T3 records a larger point estimate, $r_{\mathrm{rb}}=+0.43$, but with
$93.7\%$ of scores pinned at the maximum $1.0$ and an effective sample of
$\nsample_{\mathrm{eff}}\approx 5$ after the ceiling censors the rest. It is
tempting to read the larger $r$ as a latent structure effect that T1 lacked.
The theory says the opposite, and assigns T3 to a \emph{different} regime than
T1 --- this distinctness is the substantive content of the subsection.

T1's null lives at the information limit (\cref{thm:irrelevance}): the channel
is closed and the bounded residual is small. T3's ceiling lives instead at the
\emph{floor of the risk}, where the bounded-learner gap is inactive for a reason
unrelated to parity. When a task sits below the model's capacity frontier, the
predictor realizes (near-)Bayes behaviour in \emph{every} representation:
$\risk(\widehat f_{\rep})\to\bayesrisk(\mathrm{id})$ with
$\approxerr{\rep}\to 0$ for each C-condition simultaneously. By
\cref{cor:realizable-gap} the achievable advantage then collapses to the
estimation term alone, and with the risk already at its floor there is no
headroom in which any term --- capacity or otherwise --- can express itself: the
saturation is the \emph{realizable, infinite-resource} corner of
\cref{prop:regime}(iii), reached not by taking $\nsample\to\infty$ but by the
task being easy enough that finite resources already suffice. The bounded gap is
inactive because the frontier was never approached, which is a categorically
different statement from T1's ``inactive because the information channel is
closed.''

On top of this theoretical saturation sits a measurement one, and the two must
not be conflated. The $93.7\%$ ceiling censors the score distribution, so even
if a sub-ceiling advantage existed it could not be resolved: with
$\nsample_{\mathrm{eff}}\approx 5$ the $r_{\mathrm{rb}}=+0.43$ is an estimate
over a handful of unsaturated items and carries almost no power. The honest
reading is therefore that T3 is \emph{uninformative} about the bounded gap in
both senses at once --- the theory expects no active gap (sub-frontier
saturation) \emph{and} the instrument could not have measured one (ceiling
censoring). The larger $r$ is consistent with the theory not because the theory
predicts it but because the theory predicts the cell carries no identifying
signal in either direction; reading the point estimate as evidence of a
structure advantage would be an overclaim the theory forbids.

\idealization{The saturation argument uses \cref{cor:realizable-gap}, which
\emph{relaxes} idealization~(i) (it is a bounded-learner statement, not a Bayes
one) but assumes the base class renders $Y$ realizable in every coordinate ---
the formal content of ``below the capacity frontier.'' For a frozen LLM this
realizability is an empirical surmise about task difficulty, not a verified
property; it is the conjectural step here, and the
distinctness-from-T1 claim depends on it.}

\subsection{The T2 positive as a permitted, unidentified bounded effect}
\label{subsec:anchor-t2}

For completeness the third cell: T2 recorded the largest clean effect,
$r_{\mathrm{rb}}=+0.58$ at raw $p=0.023$, which did not survive the Holm
correction over the nine-cell confirmatory family ($m=9$; non-rejection on all
nine). The theory neither claims nor needs this as a positive. What it says is
narrow and one-directional: an effect of this magnitude is \emph{permitted} ---
it lies inside the capacity budget $\Phi$ of \cref{thm:bounded-gap} and the
two-sided envelope $|\advantage|\le\Phi$ of \cref{cor:size-gap}, so its mere
existence violates no bound --- and, were it real, it would sit in the
activation window of \cref{prop:regime}(ii), the finite-sample
capacity-separated regime the theory names as the only place a parity-preserving
effect can live. That is the whole of the retrodiction: the sign is
unconstrained by the theory, the size is bounded, and the
statistical verdict on the frozen record is non-rejection. T2 is the cell most
worth pointing a designed bounded-regime experiment at (\cref{sec:predictions}),
precisely because it is the one the theory cannot resolve from the existing
data.

\subsection{Limits of the mapping}
\label{subsec:anchor-limits}

\begin{remark}[What this section does and does not establish]
\label{rem:anchor-limits}
The three cells map cleanly onto three named regimes --- T1 to the
information-limit invariance of \cref{thm:irrelevance} (channel closed), T3 to
sub-frontier saturation via \cref{cor:realizable-gap} and
\cref{prop:regime}(iii) (no headroom, compounded by ceiling censoring), and T2
to the permitted-but-unidentified activation window of \cref{prop:regime}(ii).
The mapping's limits, stated plainly so \cref{sec:limitations} can collect them:
\begin{enumerate}
  \item[(L1)] \emph{Consistency is not confirmation.} Every claim above is of
    the form ``the record is consistent with the theory.'' A null and a ceiling
    are weak evidence; they fail to refute, which is not the same as
    corroborating. The theory earns its keep through the \emph{predictions} of
    \cref{sec:predictions}, not through this retrodiction.
  \item[(L2)] \emph{The theory does not prove any effect is zero.}
    \cref{thm:irrelevance} forces the \emph{Bayes-optimal} risks equal; it is
    silent on a bounded learner, where \cref{thm:bounded-gap} permits a nonzero
    advantage up to $\Phi$. ``No information-channel effect is expected'' is the
    exact claim; ``the effect is zero'' is not, and conflating them would be the
    central overclaim this paper takes care to prevent.
  \item[(L3)] \emph{Idealizations are not met by a frozen LLM.} The model is
    neither Bayes-optimal (idealization~(i)) nor at absolute parity
    (idealization~(ii), corrected only up to the \cref{sec:parity-defect}
    defect), and it is a single fixed predictor rather than an ERM over a class
    with a sample-size knob, so $\nsample$, $\hypclass$, and $\widehat f_{\rep}$
    enter by analogy. The retrodiction is robust to these only because it asserts
    \emph{bounds and regimes}, never point values.
  \item[(L4)] \emph{The T1/T3 distinction is partly conjectural.} That T3 sits
    below the capacity frontier (hence realizable, hence saturated) is an
    empirical surmise about task difficulty, not a theorem; if
    T3 were in fact frontier-hard, its ceiling would be measurement saturation
    without the sub-frontier reading, and only the censoring argument would
    survive.
\end{enumerate}
In one line: the theory retrodicts \emph{which regime each cell occupies} and
\emph{why no information-limit effect should appear}, and it declines ---
deliberately --- to retrodict the existence, sign, or magnitude of any
bounded-regime effect, leaving those to the falsifiable tests of
\cref{sec:predictions}.
\end{remark}
